\documentclass[10pt,a4paper]{amsart}
\usepackage{amsmath,amssymb,amsthm,graphicx}
\usepackage{verbatim}
\newenvironment{reference}{\comment}{\endcomment}
%\newenvironment{reference}{}{}
\newenvironment{slogan}{\comment}{\endcomment}
\newenvironment{history}{\comment}{\endcomment}

% For commutative diagrams we use Xy-pic
\usepackage[all]{xy}

% We use 2cell for 2-commutative diagrams.
\xyoption{2cell}
% Theorem environments.
%
\theoremstyle{plain}
\newtheorem{theorem}[subsection]{Theorem}
\newtheorem{proposition}[subsection]{Proposition}
\newtheorem{lemma}[subsection]{Lemma}

\theoremstyle{definition}
\newtheorem{definition}[subsection]{Definition}
\newtheorem{example}[subsection]{Example}
\newtheorem{exercise}[subsection]{Exercise}
\newtheorem{situation}[subsection]{Situation}

\theoremstyle{remark}
\newtheorem{remark}[subsection]{Remark}
\newtheorem{remarks}[subsection]{Remarks}

\numberwithin{equation}{subsection}

% Macros
%
\def\lim{\mathop{\mathrm{lim}}\nolimits}
\def\colim{\mathop{\mathrm{colim}}\nolimits}
\def\Spec{\mathop{\mathrm{Spec}}}
\def\Hom{\mathop{\mathrm{Hom}}\nolimits}
\def\Ext{\mathop{\mathrm{Ext}}\nolimits}
\def\SheafHom{\mathop{\mathcal{H}\!\mathit{om}}\nolimits}
\def\SheafExt{\mathop{\mathcal{E}\!\mathit{xt}}\nolimits}
\def\Sch{\mathit{Sch}}
\def\Mor{\mathop{\mathrm{Mor}}\nolimits}
\def\Ob{\mathop{\mathrm{Ob}}\nolimits}
\def\Sh{\mathop{\mathit{Sh}}\nolimits}
\def\NL{\mathop{N\!L}\nolimits}
\def\CH{\mathop{\mathrm{CH}}\nolimits}
\def\proetale{{pro\text{-}\acute{e}tale}}
\def\etale{{\acute{e}tale}}
\def\QCoh{\mathit{QCoh}}
\def\Ker{\mathop{\mathrm{Ker}}}
\def\Im{\mathop{\mathrm{Im}}}
\def\Coker{\mathop{\mathrm{Coker}}}
\def\Coim{\mathop{\mathrm{Coim}}}

% Boxtimes
%
\DeclareMathSymbol{\boxtimes}{\mathbin}{AMSa}{"02}

%
% Macros for moduli stacks/spaces
%
\def\QCohstack{\mathcal{QC}\!\mathit{oh}}
\def\Cohstack{\mathcal{C}\!\mathit{oh}}
\def\Spacesstack{\mathcal{S}\!\mathit{paces}}
\def\Quotfunctor{\mathrm{Quot}}
\def\Hilbfunctor{\mathrm{Hilb}}
\def\Curvesstack{\mathcal{C}\!\mathit{urves}}
\def\Polarizedstack{\mathcal{P}\!\mathit{olarized}}
\def\Complexesstack{\mathcal{C}\!\mathit{omplexes}}
% \Pic is the operator that assigns to X its picard group, usage \Pic(X)
% \Picardstack_{X/B} denotes the Picard stack of X over B
% \Picardfunctor_{X/B} denotes the Picard functor of X over B
\def\Pic{\mathop{\mathrm{Pic}}\nolimits}
\def\Picardstack{\mathcal{P}\!\mathit{ic}}
\def\Picardfunctor{\mathrm{Pic}}
\def\Deformationcategory{\mathcal{D}\!\mathit{ef}}
\usepackage{cleveref}
\crefname{originalreference}{original reference}{original references}
\Crefname{originalreference}{Original reference}{Original references}
\makeatletter
\newlabel{lemma-base-change-tensor}{{0A1D}{source}{}{sourcepointer-1}{}}
\newlabel{lemma-base-change-tensor@cref}{{[originalreference][0][]0A1D}{[1][0][]source}{}{}{}}
\makeatother
\crefname{originalreference}{original reference}{original references}
\Crefname{originalreference}{Original reference}{Original references}
\makeatletter
\newlabel{lemma-single-complex-base-change}{{0DJ9}{source}{}{sourcepointer-2}{}}
\newlabel{lemma-single-complex-base-change@cref}{{[originalreference][0][]0DJ9}{[1][0][]source}{}{}{}}
\makeatother
\crefname{originalreference}{original reference}{original references}
\Crefname{originalreference}{Original reference}{Original references}
\makeatletter
\newlabel{lemma-single-complex-base-change-condition}{{0DJ8}{source}{}{sourcepointer-3}{}}
\newlabel{lemma-single-complex-base-change-condition@cref}{{[originalreference][0][]0DJ8}{[1][0][]source}{}{}{}}
\makeatother
\crefname{originalreference}{original reference}{original references}
\Crefname{originalreference}{Original reference}{Original references}
\makeatletter
\newlabel{lemma-single-complex-base-change-condition-inherited}{{0DJA}{source}{}{sourcepointer-4}{}}
\newlabel{lemma-single-complex-base-change-condition-inherited@cref}{{[originalreference][0][]0DJA}{[1][0][]source}{}{}{}}
\makeatother
\title{Base change for internal Hom from perfect or pseudo-coherent objects}
\author{The Stacks Project Authors}
\date{stacks-project:spaces-perfect.tex:local-snapshot-20261012; GNU FDL 1.2 or later}
\begin{document}
\maketitle
\noindent\textit{Editorial note.} All delivered mathematical text is the author's own, copied from the pinned Stacks Project chapter file spaces-perfect.tex; no statement or proof step was rewritten, reordered or removed.\par\medskip
\noindent\textit{Editorial note.} The counted proof is the author's complete argument: the proof states the precise base change assertion, reduces via the single-complex base change condition and its inheritance, and refers the stalk verification to the argument of the tensor base change lemma in the same chapter\par\medskip
\noindent\textit{Editorial note.} The label lemma-base-change-RHom names the counted statement, and the proof environment immediately following it in the source is delivered verbatim as its proof.\par\medskip
\noindent\textit{Editorial note.} The counted statement is carried by the author's own lemma environment in the chapter source.\par\medskip
\noindent\textit{Allowed standard prerequisites (editorial).}\begin{itemize}
\item the single-complex base change lemma, its condition, and the inheritance of the condition (lemma-single-complex-base-change, lemma-single-complex-base-change-condition, lemma-single-complex-base-change-condition-inherited of the same chapter)
\item the stalk verification that for a complex G flat over Y the canonical map L(g')\textasciicircum{}*G -> (g')\textasciicircum{}*G satisfies the base change condition, as established in the proof of lemma-base-change-tensor of the same chapter
\end{itemize}\medskip
\ifdefined\hypertarget\hypertarget{sourcepointer-1}{}\fi
\noindent\textit{Source reference.} Original reference 0A1D: sources/stacks-spaces-perfect/pinned-original-source.tex, source lines 5625--5625. Prerequisite pointer invoked at proof line(s) 5720; it names a standard fact of the same chapter that the counted argument uses, and every step of the counted proof built on it is retained verbatim in the delivered body.\par\medskip
\ifdefined\hypertarget\hypertarget{sourcepointer-2}{}\fi
\noindent\textit{Source reference.} Original reference 0DJ9: sources/stacks-spaces-perfect/pinned-original-source.tex, source lines 5488--5488. Prerequisite pointer invoked at proof line(s) 5712; it names a standard fact of the same chapter that the counted argument uses, and every step of the counted proof built on it is retained verbatim in the delivered body.\par\medskip
\ifdefined\hypertarget\hypertarget{sourcepointer-3}{}\fi
\noindent\textit{Source reference.} Original reference 0DJ8: sources/stacks-spaces-perfect/pinned-original-source.tex, source lines 5365--5365. Prerequisite pointer invoked at proof line(s) 5719; it names a standard fact of the same chapter that the counted argument uses, and every step of the counted proof built on it is retained verbatim in the delivered body.\par\medskip
\ifdefined\hypertarget\hypertarget{sourcepointer-4}{}\fi
\noindent\textit{Source reference.} Original reference 0DJA: sources/stacks-spaces-perfect/pinned-original-source.tex, source lines 5567--5567. Prerequisite pointer invoked at proof line(s) 5713; it names a standard fact of the same chapter that the counted argument uses, and every step of the counted proof built on it is retained verbatim in the delivered body.\par\medskip
\begin{lemma}
\label{lemma-base-change-RHom}
Let $S$ be a scheme. Let $f : X \to Y$ be a quasi-compact and
quasi-separated morphism of algebraic spaces over $S$. Let $E$
be an object of $D(\mathcal{O}_X)$. Let $\mathcal{G}^\bullet$
be a complex of quasi-coherent $\mathcal{O}_X$-modules. If
\begin{enumerate}
\item $E$ is perfect, $\mathcal{G}^\bullet$ is a bounded above,
and $\mathcal{G}^n$ is flat over $Y$, or
\item $E$ is pseudo-coherent, $\mathcal{G}^\bullet$ is bounded,
and $\mathcal{G}^n$ is flat over $Y$,
\end{enumerate}
then formation of
$$
Rf_*R\SheafHom(E, \mathcal{G}^\bullet)
$$
commutes with arbitrary base change (see proof for precise statement).
\end{lemma}

\begin{proof}
The statement means the following. Let $g : Y' \to Y$ be a morphism of
algebraic spaces and consider the base change diagram
$$
\xymatrix{
X' \ar[r]_h \ar[d]_{f'} &
X \ar[d]^f \\
Y' \ar[r]^g &
Y
}
$$
in other words $X' = Y' \times_Y X$. The lemma asserts that
$$
Lg^*Rf_*R\SheafHom(E, \mathcal{G}^\bullet)
\longrightarrow
R(f')_*R\SheafHom(L(g')^*E, (g')^*\mathcal{G}^\bullet)
$$
is an isomorphism. Observe that on the right hand side we do {\bf not}
use the derived pullback on $\mathcal{G}^\bullet$. To prove this, we apply
Lemmas \ref{lemma-single-complex-base-change} and
\ref{lemma-single-complex-base-change-condition-inherited} to see that it
suffices to prove the canonical map
$$
L(g')^*\mathcal{G}^\bullet \to (g')^*\mathcal{G}^\bullet
$$
satisfies the equivalent conditions of
Lemma \ref{lemma-single-complex-base-change-condition}.
This was shown in the proof of Lemma \ref{lemma-base-change-tensor}.
\end{proof}

\end{document}
